\documentclass{article}

\usepackage{../../Style}

\title{Group Cohomology}
\author{Shing Tak Lam}
\date{Lent 2023\footnote{Based on lectures by Chris Brookes. Last updated \today.}}

\makeindex

\begin{document}

\maketitle

\tableofcontents

\lec

\section{Definitions and resolutions}

Let \(G\) be a group.

\begin{definition}
    [integral group ring]
    The \emph{integral group ring}\index{integral group ring} \(\ZG\) has elements being formal sums
    \[\sum n_g g\]
    for \(n_g \in \Z, g \in G\), and only finitely many \(n_g\) are non-zero. This is a free abelian group under addition, i.e.
    \[\left(\sum m_g g\right) + \left(\sum n_g g\right) = \sum (m_g + n_g) g\]
    and we have multiplication
    \[\left(\sum m_h h\right)\left(\sum n_k k\right) = \sum_g \left(\sum_{hk = g}m_hn_k\right) g\]
\end{definition}

This is the associative ring underlying the integral representation theory of \(G\). We will write \(1 = 1e\) for the multiplicative identity in \(G\).

\begin{definition}
    [module]
    A \emph{(left) \(\ZG\)-module}\index{ZG module@\(\ZG\) module} \(M\) is an abelian group under addition, and a map
    \begin{align*}
        \ZG \times M &\to M \\
        (r, m) &\mapsto rm
    \end{align*}
    such that
    \begin{enumerate}
        \item \(r(m + n) = rm + rn\),
        \item \((r_1 + r_2)m = r_1m + r_2 m\),
        \item \((r_1r_2)m = r_1(r_2m)\),
        \item \(1m = m\)
    \end{enumerate}
\end{definition}

A module is \emph{trivial} if \(gm = m\) for all \(g \in G\), \(m \in M\). We say \emph{the trivial module}\index{trivial \(\ZG\) module} is \(\Z\) with trivial action \(gm = m\). A \emph{free \(\ZG\)-module}\index{free \(\ZG\) module} on \(X\) is formal sums
\[\sum r_x x\]
for \(r_x \in \ZG, x \in X\), and finitely many \(r_x\) non-zero. We will write \(\ZG\{X\}\) for this.

\emph{Submodules} and \emph{quotient modules} are defined as we would expect.

\begin{definition}
    [\(\ZG\)-map]
    An \emph{\(\ZG\)-map (or morphism)}\index{\(\ZG\)-map} \(\alpha : M_1 \to M_2\) is a map of abelian groups, with
    \[\alpha(rm_1) = r\alpha(m_1)\]
    for \(r \in \ZG, m_1 \in M_1\).
\end{definition}

\begin{example}
    The \emph{augmentation map}\index{augmentation map} is
    \begin{align*}
        \epsilon : \ZG &\to \Z \\
        \sum n_g g &\mapsto \sum n_g
    \end{align*}
    where we treat \(\ZG\) as a left \(\ZG\)-module, and \(\Z\) as the trivial module. This is a \(\ZG\) map, and it is also a ring homomorphism.
\end{example}

\begin{notation}
    Write \(\Hom_G(M, N)\)\index{HomGMN@\(\Hom_G(M, N)\)} for the set of \(\ZG\)-maps from \(M \to N\), with pointwise addition.
\end{notation}

\begin{example}
    If we regard \(\ZG\) as a left \(\ZG\)-module, then
    \[\Hom_G(\ZG, M) \cong M\]
    for any left \(\ZG\)-module \(M\), by sending \(\varphi\) to \(\varphi(1)\). This is an isomorphism since
    \[\varphi(r) = \varphi(r1) = r\varphi(1)\]
    Note that \(\Hom_G(\ZG, M)\) is a left \(\ZG\)-module, with
    \[(s\phi)(r) = \phi(rs)\]
    In particular,
    \[\Hom_G(\ZG, \ZG) \cong \ZG\]
    With this, \(\varphi \mapsto \varphi(1)\) corresponds to multiplication on the \emph{right} by \(\varphi(1)\).
\end{example}

Note that \(G\) may not be abelian, so the distinction between left and right modules matters.

\begin{definition}
    If \(f : M_1 \to M_2\) is a \(\ZG\)-map, then the dual map is
    \[f^* : \Hom_G(M_2, N) \to \Hom_G(M_1, N)\]
    for any \(\ZG\)-module \(N\). This is given by
    \[f^*(\varphi) = \varphi \circ f\]

    Similarly, if \(f : N_1 \to N_2\), then we have
    \[f_* : \Hom_G(M, N_1) \to \Hom(M, N_2)\]
    with \(f_*(\varphi) = f \circ \varphi\).
\end{definition}

These are maps of abelian groups, since in general \(\Hom_G(M, N)\) doesn't have to be a \(\ZG\)-module, as we don't have a right action.

\begin{example}
    [Prototypical example] Let \(G = \inner{t}\) be an infinite cyclic group, and consider the graph \(\Gamma\) with vertices \(\{v_i\}_{i \in \Z}\), edges \(v_i \bij v_{i+1}\). Then \(G\) acts on the set \(V\) of vertices, by \(t v_i = v_{i+1}\), and also on the set \(E\) of edges. The action is transitive in both cases. The formal integral sums \(\Z V\) and \(\Z E\) are \(\ZG\)-modules, and are free. Fix the edge \(e : v_0 \bij v_1\). Then we have \(\ZG\)-maps
    \begin{align*}
        \dd : \Z E &\to \Z V \\
        e &\mapsto v_1 - v_0
    \end{align*}
    and also
    \begin{align*}
        \Z V &\to \Z \\
        v_0 &\mapsto 1
    \end{align*}
    This correspond to the augmentation map.
\end{example}

\begin{definition}
    [chain complex, exact, homology]
    A \emph{chain complex}\index{chain complex} of \(\ZG\)-modules is a sequence
    % https://q.uiver.app/#q=WzAsNCxbMCwwLCJNX3MiXSxbMSwwLCJNX3tzLTF9Il0sWzIsMCwiXFxjZG90cyJdLFszLDAsIk1fdCJdLFswLDEsImRfcyJdLFsxLDIsImRfe3MtMX0iXSxbMiwzLCJkX3t0KzF9Il1d
    \[\begin{tikzcd}
        {M_s} & {M_{s-1}} & \cdots & {M_t}
        \arrow["{\dd_s}", from=1-1, to=1-2]
        \arrow["{\dd_{s-1}}", from=1-2, to=1-3]
        \arrow["{\dd_{t+1}}", from=1-3, to=1-4]
    \end{tikzcd}\]
    of \(\ZG\)-modules and maps, such that for \(t < n < s\), \(\dd_n \circ \dd_{n+1} = 0\). We will write
    \[M_\bullet = (M_n, \dd_n)_{t \le n \le s}\]
\end{definition}
We say that \(M_\bullet\) is \emph{exact}\index{exact} at \(M_n\) if \(\im(\dd_{n+1}) = \ker(\dd n)\). The sequence is exact if if it exact at all \(M_n\) with \(t < n < s\).

The \emph{homology}\index{homology} of \(M_\bullet\) is given by
\begin{align*}
    \HH_s(M_\bullet) &= \ker(\dd s) \\
    \HH_n(M_\bullet) &= \frac{\ker(\dd_n)}{\im(\dd_{n+1})} \text{ for }t < n < s \\
    \HH_t(M_\bullet) &= \coker(\dd_{t+1})
\end{align*}

A \emph{short exact sequence}\index{short exact sequence} is a chain complex of the form
% https://q.uiver.app/#q=WzAsNSxbMCwwLCIwIl0sWzEsMCwiTV8xIl0sWzIsMCwiTV8yIl0sWzMsMCwiTV8zIl0sWzQsMCwiMCJdLFswLDFdLFsxLDIsIlxcYWxwaGEiXSxbMiwzLCJcXGJldGEiXSxbMyw0XV0=
\[\begin{tikzcd}
	0 & {M_1} & {M_2} & {M_3} & 0
	\arrow[from=1-1, to=1-2]
	\arrow["\alpha", from=1-2, to=1-3]
	\arrow["\beta", from=1-3, to=1-4]
	\arrow[from=1-4, to=1-5]
\end{tikzcd}\]
In our prototypical example, we have a short exact sequence
% https://q.uiver.app/#q=WzAsNSxbMCwwLCIwIl0sWzEsMCwiXFxaIEUiXSxbMiwwLCJcXFogViJdLFszLDAsIlxcWiAiXSxbNCwwLCIwIl0sWzAsMV0sWzEsMl0sWzIsM10sWzMsNF1d
\[\begin{tikzcd}
	0 & {\Z E} & {\Z V} & {\Z } & 0
	\arrow[from=1-1, to=1-2]
	\arrow[from=1-2, to=1-3]
	\arrow[from=1-3, to=1-4]
	\arrow[from=1-4, to=1-5]
\end{tikzcd}\]
which corresponds to
% https://q.uiver.app/#q=WzAsNSxbMCwwLCIwIl0sWzMsMCwiXFxaICJdLFs0LDAsIjAiXSxbMSwwLCJcXFogRyJdLFsyLDAsIlxcWiBHIl0sWzEsMl0sWzAsM10sWzMsNF0sWzQsMV1d
\[\begin{tikzcd}
	0 & {\ZG} & {\ZG} & {\Z } & 0
	\arrow[from=1-4, to=1-5]
	\arrow[from=1-1, to=1-2]
	\arrow[from=1-2, to=1-3]
	\arrow[from=1-3, to=1-4]
\end{tikzcd}\]
where the map \(\ZG \to \ZG\) is multiplication on the right by \(t - 1\).

\begin{definition}
    [projective]
    A \(\ZG\)-module \(P\) is \emph{projective}\index{projective} if for every surjective \(\ZG\)-map \(\alpha : M_1 \to M_2\), and every \(\ZG\)-map \(\beta : P \to M_2\), there exists \(\overline\beta\) such that
    % https://q.uiver.app/#q=WzAsNCxbMCwxLCJNXzEiXSxbMSwxLCJNXzIiXSxbMiwxLCIwIl0sWzEsMCwiUCJdLFswLDEsIlxcYWxwaGEiLDJdLFsxLDJdLFszLDEsIlxcYmV0YSJdLFswLDMsIlxcb3ZlcmxpbmVcXGJldGEiXV0=
    \[\begin{tikzcd}
        & P \\
        {M_1} & {M_2} & 0
        \arrow["\alpha"', from=2-1, to=2-2]
        \arrow[from=2-2, to=2-3]
        \arrow["\beta", from=1-2, to=2-2]
        \arrow["\overline\beta", from=2-1, to=1-2]
    \end{tikzcd}\]
    commutes.
\end{definition}
If we have any short exact sequence
% https://q.uiver.app/#q=WzAsNSxbMCwwLCIwIl0sWzEsMCwiTiJdLFsyLDAsIk1fMSJdLFszLDAsIk1fMiJdLFs0LDAsIjAiXSxbMCwxXSxbMSwyLCJmIl0sWzIsMywiXFxhbHBoYSJdLFszLDRdXQ==
\[\begin{tikzcd}
	0 & N & {M_1} & {M_2} & 0
	\arrow[from=1-1, to=1-2]
	\arrow["f", from=1-2, to=1-3]
	\arrow["\alpha", from=1-3, to=1-4]
	\arrow[from=1-4, to=1-5]
\end{tikzcd}\]
Consider
% https://q.uiver.app/#q=WzAsNSxbMCwwLCIwIl0sWzEsMCwiXFxIb21fRyhQLCBOKSJdLFsyLDAsIlxcSG9tX0coUCwgTV8xKSJdLFszLDAsIlxcSG9tX0coUCwgTV8yKSJdLFs0LDAsIjAiXSxbMCwxXSxbMSwyLCJmXyoiXSxbMiwzLCJcXGFscGhhXyoiXSxbMyw0XV0=
\[\begin{tikzcd}
	0 & {\Hom_G(P, N)} & {\Hom_G(P, M_1)} & {\Hom_G(P, M_2)} & 0
	\arrow[from=1-1, to=1-2]
	\arrow["{f_*}", from=1-2, to=1-3]
	\arrow["{\alpha_*}", from=1-3, to=1-4]
	\arrow[from=1-4, to=1-5]
\end{tikzcd} \tag{\(*\)}\]
Then \(P\) is projective if and only if \((*)\) is exact. Note that in \((*)\), we always have exactness except at \(\Hom_G(P, M_2)\).

\lec

\begin{lemma}
    Free modules are projective.
\end{lemma}

\begin{proof}
    Let \(\alpha : M_1 \to M_2\) be a surjective \(\ZG\)-map, \(\beta : \ZG\{X\} \to M_2\). For each \(x \in X\), since \(\alpha\) is surjective, there exists \(m_x \in M_1\) such that \(\alpha(m_x) = \beta(x)\). We can then define \(\overline\beta : \ZG\{X\} \to M_1\), by
    \[\sum r_x \cdot x \mapsto \sum r_x \cdot m_x\]
\end{proof}

\begin{definition}
    A \emph{projective (resp. free) resolution of the trivial module \(\Z\)} is an exact sequence
    % https://q.uiver.app/#q=WzAsNSxbMCwwLCJcXGNkb3RzIl0sWzEsMCwiUF8xIl0sWzIsMCwiUF8wIl0sWzMsMCwiXFxaIl0sWzQsMCwiMCJdLFswLDEsIlxcZGRfMiJdLFsxLDIsIlxcZGRfISJdLFsyLDMsIlxcZGRfMCJdLFszLDRdXQ==
    \[\begin{tikzcd}
        \cdots & {P_1} & {P_0} & \Z & 0
        \arrow["{\dd_2}", from=1-1, to=1-2]
        \arrow["{\dd_!}", from=1-2, to=1-3]
        \arrow["{\dd_0}", from=1-3, to=1-4]
        \arrow[from=1-4, to=1-5]
    \end{tikzcd}\]
    with all the \(P_i\) projective (resp. free).
\end{definition}

\begin{example}
    When \(G = \inner{t}\) is infinite cyclic, we have the free resolution
    % https://q.uiver.app/#q=WzAsNSxbMCwwLCIwIl0sWzEsMCwiXFxaRyJdLFsyLDAsIlxcWkciXSxbNCwwLCIwIl0sWzMsMCwiXFxaIl0sWzAsMV0sWzEsMiwiXFxjZG90KHQtMSkiXSxbMiw0LCJcXGVwc2lsb24iXSxbNCwzXV0=
\[\begin{tikzcd}
	0 & \ZG & \ZG & \Z & 0
	\arrow[from=1-1, to=1-2]
	\arrow["{\cdot(t-1)}", from=1-2, to=1-3]
	\arrow["\epsilon", from=1-3, to=1-4]
	\arrow[from=1-4, to=1-5]
\end{tikzcd}\]
\end{example}

\begin{example}
    If \(G = \inner{t}\) cyclic of order \(n\). Then
    % https://q.uiver.app/#q=WzAsNyxbNiwwLCIwIl0sWzUsMCwiXFxaIl0sWzQsMCwiXFxaRyJdLFszLDAsIlxcWkciXSxbMiwwLCJcXFpHIl0sWzEsMCwiXFxaRyJdLFswLDAsIlxcY2RvdHMiXSxbMiwxLCJcXGVwc2lsb24iXSxbMSwwXSxbMywyLCJcXGFscGhhIl0sWzQsMywiXFxiZXRhIl0sWzUsNCwiXFxhbHBoYSJdLFs2LDVdXQ==
\[\begin{tikzcd}
	\cdots & \ZG & \ZG & \ZG & \ZG & \Z & 0
	\arrow["\epsilon", from=1-5, to=1-6]
	\arrow[from=1-6, to=1-7]
	\arrow["\alpha", from=1-4, to=1-5]
	\arrow["\beta", from=1-3, to=1-4]
	\arrow["\alpha", from=1-2, to=1-3]
	\arrow[from=1-1, to=1-2]
\end{tikzcd}\]
    is a free resolution, where \(\alpha(x) = x(t-1)\) and \(\beta(x) = x(1 + t + \dots + t^{n-1})\).
\end{example}

From algebraic topology, if we have a connected simplicial complex \(X\), with \(\pi_1(X) = G\), and the universal cover \(\widetilde X\) is contractible, then we have a free resolution of \(\Z\). The point is that the simplicial complex \(X\) contains a lot of information about \(G\). We are trying to replicate this algebraically.

For calculation purposes, we're interested in \emph{small resolutions}. For example, where the free modules have small rank. But for theory development purposes, we're wanting general constructions, and such resolutions tend to be large.

\begin{definition}
    \(G\) is of type \(\FP_n\) if \(\Z\) has a projective resolution
    % https://q.uiver.app/#q=WzAsNixbMCwwLCJcXGNkb3RzIl0sWzIsMCwiUF8xIl0sWzEsMCwiUF8yIl0sWzMsMCwiUF8wIl0sWzQsMCwiXFxaIl0sWzUsMCwiMCJdLFswLDJdLFsyLDEsIlxcZGRfMiJdLFsxLDMsIlxcZGRfMSJdLFszLDQsIlxcZGRfMCJdLFs0LDVdXQ==
    \[\begin{tikzcd}
        \cdots & {P_2} & {P_1} & {P_0} & \Z & 0
        \arrow[from=1-1, to=1-2]
        \arrow["{\dd_2}", from=1-2, to=1-3]
        \arrow["{\dd_1}", from=1-3, to=1-4]
        \arrow["{\dd_0}", from=1-4, to=1-5]
        \arrow[from=1-5, to=1-6]
    \end{tikzcd}\]
    so that \(P_n, \dots, P_0\) are finitely generated. 
    
    \(G\) is of type \(\FP_n\) is \(\Z\) has a projective resolution with all \(P_n\) finitely generated. 
    
    \(G\) is of type \(\FP\) if there exists a projective resolution of \(\Z\) with finite length and all \(P_n\) are finitely generated. That is, \(P_s = 0\) for \(s\) sufficiently large.
\end{definition}

\begin{example}
    \(G = \inner{t}\) infinite cyclic is of type \(\FP\), \(G = C_n = \inner{t}\) finite cyclic is of type \(\FP_\infty\).
\end{example}

These are to be regarded as finiteness conditions on the group. The topological version of \(\FP_n\) (which is called \(\mathrm{F}_n\)) would be asking for \(X\) with \(\pi_1(X) = G\), having a finite \(n\)-skeleton. Note that these two conditions are not equivalent.

Now lets meet some general constructions. If we have a partial projective resolution
% https://q.uiver.app/#q=WzAsNyxbMCwwLCJQX3MiXSxbMSwwLCJQX3tzLTF9Il0sWzIsMCwiXFxjZG90cyJdLFszLDAsIlBfMSJdLFs0LDAsIlBfMCJdLFs1LDAsIlxcWiJdLFs2LDAsIjAiXSxbMCwxLCJcXGRkX3MiXSxbMiwzXSxbMSwyXSxbMyw0LCJcXGRkXzEiXSxbNCw1LCJcXGRkXzAiXSxbNSw2XV0=
\[\begin{tikzcd}
	{P_s} & {P_{s-1}} & \cdots & {P_1} & {P_0} & \Z & 0
	\arrow["{\dd_s}", from=1-1, to=1-2]
	\arrow[from=1-3, to=1-4]
	\arrow[from=1-2, to=1-3]
	\arrow["{\dd_1}", from=1-4, to=1-5]
	\arrow["{\dd_0}", from=1-5, to=1-6]
	\arrow[from=1-6, to=1-7]
\end{tikzcd}\]
Then we can make it longer by setting
\[P_{s+1} = \ZG\{X_{s+1}\}\]
with \(X_{s+1} = \ker(\dd_s)\), and set
\[\dd_{s+1} : \sum r_x x \mapsto \sum r_x x\]
Then
% https://q.uiver.app/#q=WzAsOCxbMSwwLCJQX3MiXSxbMiwwLCJQX3tzLTF9Il0sWzMsMCwiXFxjZG90cyJdLFs0LDAsIlBfMSJdLFs1LDAsIlBfMCJdLFs2LDAsIlxcWiJdLFs3LDAsIjAiXSxbMCwwLCJQX3tzKzF9Il0sWzAsMSwiXFxkZF9zIl0sWzIsM10sWzEsMl0sWzMsNCwiXFxkZF8xIl0sWzQsNSwiXFxkZF8wIl0sWzUsNl0sWzcsMCwiXFxkZF97cysxfSJdXQ==
\[\begin{tikzcd}
	{P_{s+1}} & {P_s} & {P_{s-1}} & \cdots & {P_1} & {P_0} & \Z & 0
	\arrow["{\dd_s}", from=1-2, to=1-3]
	\arrow[from=1-4, to=1-5]
	\arrow[from=1-3, to=1-4]
	\arrow["{\dd_1}", from=1-5, to=1-6]
	\arrow["{\dd_0}", from=1-6, to=1-7]
	\arrow[from=1-7, to=1-8]
	\arrow["{\dd_{s+1}}", from=1-1, to=1-2]
\end{tikzcd}\]
is a longer partial projective resolution. To make \(P_{s+1}\) smaller, we could take \(X_{s+1}\) to be a \(\ZG\)-generating set for \(\ker(\dd_s)\). This is particularly useful if \(\ker(\dd_s)\) is finitely generated.

Continuing this, we can get a (full) projective resolution.

\begin{definition}
    The \emph{standard (or bar) resolution of \(\Z\)}, for any group \(G\), is: Let
    \[G^{(n)} = \{[g_1\vert\dots\vert g_n] \mid g_i \in G\}\]
    be the set of symbols. In particular, \(G^{(0)} = \{[]\}\) Let
    \[F_n = \ZG\{G^{(n)}\}\]
    be the free module on \(G^{(n)}\). The differential is given by
    \begin{align*}
        \dd_n[g_1\vert\dots\vert g_n] &= g_1[g_2\vert\dots\vert g_n] \\
        &- [g_1g_2\vert g_3\vert\dots\vert g_n] \\
        &+ [g_1\vert g_2g_3\vert\dots\vert g_n] + \dots \\
        &+ (-1)^{n-1}[g_1\vert\dots\vert g_{n-1}g_n] \\
        &+ [g_1\vert\dots\vert g_{n-1}]
    \end{align*}
    With this, \(\dd_{n-1} \dd_n = 0\), and we have a chain complex
    % https://q.uiver.app/#q=WzAsNyxbMCwwLCJcXGNkb3RzIl0sWzEsMCwiRl9uIl0sWzIsMCwiRl97bi0xfSJdLFszLDAsIlxcY2RvdHMiXSxbNCwwLCJGXzAiXSxbNSwwLCJcXFoiXSxbNiwwLCIwIl0sWzQsNSwiW11cXG1hcHN0byAxIl0sWzAsMV0sWzEsMl0sWzIsM10sWzMsNF0sWzUsNl1d
    \[\begin{tikzcd}
        \cdots & {F_n} & {F_{n-1}} & \cdots & {F_0} & \Z & 0
        \arrow["{[]\mapsto 1}", from=1-5, to=1-6]
        \arrow[from=1-1, to=1-2]
        \arrow[from=1-2, to=1-3]
        \arrow[from=1-3, to=1-4]
        \arrow[from=1-4, to=1-5]
        \arrow[from=1-6, to=1-7]
    \end{tikzcd}\]
\end{definition}

\begin{remark}
    The bar resolution corresponds to the standard resolution in algebraic topology. Consider \(n+1\) tuples \(G^{n+1}\), and form the free abelian group \(\Z G^{n+1}\) under addition. \(G\) acts on \(G^{n+1}\) diagonally, so
    \[g\vdot (h_0, \dots, h_n) = (gh_0, \dots, gh_n)\]
    So \(\Z G^{n+1}\) is a free \(\ZG\)-module, on the basis of \(n+1\) tuples with \(g_0 = 1\). Then we have a correspondence
    \[[g_1\vert\dots\vert g_n] \mapsto (1, g_1, g_1g_2, \dots, g_1\cdots g_n)\]
    Note that removing the first entry gives \(g_1(1, g_2, g_2g_3, \dots, g_2\cdots g_n)\), and removing the seconc entry in the tuple gives \((1, g_1g_2, g_1g_2g_3, \dots, g_1\cdots g_n)\).
\end{remark}

\begin{lemma}
    The bar resolution is exact.
\end{lemma}

\begin{proof}
    We'll just consider the \(\dd_n\) as maps of abelian groups. \(F_n\) has basis \(G \times G^{(n)}\) as a free abelian group, and \(G \times G^{(n)}\) is the set
    \[\{g_0[g_1\vert\dots\vert g_n] \mid g_i \in G\}\]
    Define \(\Z\)-linear maps \(s_n : F_n \to F_{n+1}\), such that
    \[\id_{F_n} = \dd_{n+1}s_n + s_{n-1}\dd_n \tag{\(*\)}\]
    given by
    \[s_n(g_0[g_1\vert\dots\vert g_n]) = [g_0\vert\dots\vert g_n]\]
    We can verify that \((*)\) holds on the basis, and so if \(x \in \ker(\dd_{n+1})\), then
    \[x = \id(x) = \dd_{n+1}s_n(x) + s_{n-1}\dd_n(x) = \dd_{n+1}s_n(x) \in \im(\dd_{n+1})\]
    That is, \(s_n\) is a chain homotopy from the identity to zero.
\end{proof}

\begin{corollary}
    A finite group is of type \(\FP_\infty\).
\end{corollary}

\begin{proof}
    The bar resolution gives such a resolution.
\end{proof}

\lec

\begin{definition}
    Consider a projective resolution
    % https://q.uiver.app/#q=WzAsOCxbMCwwLCJcXGNkb3RzIl0sWzEsMCwiUF97bisxfSJdLFsyLDAsIlBfbiJdLFszLDAsIlxcY2RvdHMiXSxbNCwwLCJQXzEiXSxbNSwwLCJQXzAiXSxbNiwwLCJcXFoiXSxbNywwLCIwIl0sWzAsMV0sWzEsMl0sWzIsMywiZF9uIl0sWzMsNF0sWzQsNV0sWzUsNl0sWzYsN11d
\[\begin{tikzcd}
	\cdots & {P_{n+1}} & {P_n} & \cdots & {P_1} & {P_0} & \Z & 0
	\arrow[from=1-1, to=1-2]
	\arrow[from=1-2, to=1-3]
	\arrow["{d_n}", from=1-3, to=1-4]
	\arrow[from=1-4, to=1-5]
	\arrow[from=1-5, to=1-6]
	\arrow[from=1-6, to=1-7]
	\arrow[from=1-7, to=1-8]
\end{tikzcd}\]
    of \(\Z\) by \(\ZG\)-modules. Let \(M\) be a (left) \(\ZG\)-module. Apply the functor \(\Hom_G(\cdot, M)\), we get the \emph{dual sequence}
    % https://q.uiver.app/#q=WzAsNSxbNCwwLCJcXEhvbShQXzAsIE0pIl0sWzMsMCwiXFxjZG90cyJdLFsyLDAsIlxcSG9tKFBfbiwgTSkiXSxbMSwwLCJcXEhvbShQX3tuKzF9LCBNKSJdLFswLDAsIlxcY2RvdHMiXSxbMyw0XSxbMiwzLCJkXntuKzF9IiwyXSxbMSwyXSxbMCwxXV0=
\[\begin{tikzcd}
	\cdots & {\Hom(P_{n+1}, M)} & {\Hom(P_n, M)} & \cdots & {\Hom(P_0, M)}
	\arrow[from=1-2, to=1-1]
	\arrow["{\dd^{n+1}}"', from=1-3, to=1-2]
	\arrow[from=1-4, to=1-3]
	\arrow[from=1-5, to=1-4]
\end{tikzcd} \tag{\(\dagger\)}\]
    where \(\dd^n = \dd_n^*\). Then the \emph{\(n\)-th cohomology group \(\HH^n(G, M)\) with coefficients in \(M\) is:}
    \begin{align*}
        \HH^n(G, M) &= \frac{\ker(\dd^{n+1})}{\im(\dd^n)} \text{ for }n \ge 1 \\
        \HH^0(G, M) &= \ker(\dd^1)
    \end{align*}
\end{definition}

\begin{remark}
    We have dropped the \(\Z\)-term in \((\dagger)\).

    Also, these cohomology groups are the homology groups of the chain complex
    \[C_n = \Hom_G(P_{-n}, M)\]
    with \(-\infty < n \le 0\).
\end{remark}

Later, we will see that these cohomology groups are independent of the choice of resolution. 

\begin{example}
    In our prototypical example, \(G = \inner{t}\) infinite cyclic, we have a projective resolution
    % https://q.uiver.app/#q=WzAsNSxbMCwwLCIwIl0sWzEsMCwiXFxaRyJdLFsyLDAsIlxcWkciXSxbMywwLCJcXFoiXSxbNCwwLCIwIl0sWzAsMV0sWzEsMiwiXFxkZCJdLFsyLDNdLFszLDRdXQ==
\[\begin{tikzcd}
	0 & \ZG & \ZG & \Z & 0
	\arrow[from=1-1, to=1-2]
	\arrow["\dd", from=1-2, to=1-3]
	\arrow[from=1-3, to=1-4]
	\arrow[from=1-4, to=1-5]
\end{tikzcd}\]
    where \(\dd(x) = x(t-1)\). For \(\phi \in \Hom_G(\ZG, M)\), \(x \in \ZG\),
    \[\dd^1(\phi)(x) = \phi(\dd_1(x)) = \phi(x(t-1))\]
    Moreover, recall \(\Hom_G(\ZG, M) \cong M\), where \(i(\phi) = \phi(1)\). In particular,
    \[\dd^1(\phi) \mapsto \dd^1\phi(1) = \phi(t - 1) = (t - 1)\phi(1) = (t - 1)i(\phi)\]
    In this case, we have
    % https://q.uiver.app/#q=WzAsMyxbMCwwLCIwIl0sWzEsMCwiTSJdLFsyLDAsIk0iXSxbMiwxLCIodC0xKVxcY2RvdCJdLFsxLDBdXQ==
\[\begin{tikzcd}
	0 & M & M
	\arrow["{(t-1)\cdot}", from=1-3, to=1-2]
	\arrow[from=1-2, to=1-1]
\end{tikzcd}\]
    and so
    \begin{align*}
        \HH^0(G, M) &= \{m \in M \mid tm = m\} = M^G \\
        \HH^1(G, M) &= \frac{M}{(t-1)M} = M_G
    \end{align*}
    and \(\HH^n(G, M) = 0\) otherwise. Here, \(M^G\) is the group of \emph{invariants}, the largest submodule with trivial \(G\)-action. Conversely, \(M_G\) is the \emph{coinvariants}, the largest quotient module with trivial \(G\)-action.
\end{example}

\begin{remark}
    \(\HH^0(M, G) = M^G\) is generally true. But \(\HH^1(M, G) = M_G\) is special behaviour of \(G = \inner{t}\).

    For any group of type \(\FP\), we have that \(\HH^n(G, M) = 0\) for large enough \(n\).
\end{remark}

\begin{definition}
    \(G\) is of \emph{cohomological dimension \(m\) over \(\Z\)} if there exists an \(\ZG\)-module \(M\), with 
    \[\HH^m(G, M) \ne 0\]
    and \(\HH^n(G, M') = 0\) for any \(\ZG\)-module \(M'\), and \(n > m\). 
\end{definition}

Note that for all \(G\), \(\HH^0(G, \Z) = \Z \ne 0\).

\begin{example}
    \(G = \inner{t}\) infinite cyclic is of cohomological dimension \(1\) (over \(\Z\)).

    An exercise is to show that if \(G\) is a free group of finite rank, then it is also of cohomological dimension \(1\). To see this, the Cayley graph is a tree and we can construct a resultion using this.
\end{example}

\begin{remark}
    The converse is true. A finitely generated group of cohomological dimension \(1\) is free. In fact, this is true in general. See [Stallings, 1968], {Swan, 1969}.
\end{remark}

Now consider the bar resolution in our definition of cohomology. Note that
\[\Hom_G(\ZG\{G^{(n)}\}, M) \cong C^n(G, M) = \left\{\text{functions }G^{(n)} \to M\right\}\]
This is the same as function \(G^n \to M\). We also have that
\[C^0(G, M) = \{\text{functions }[] \to M\} = M\]

\begin{definition}
    The group of \emph{\(n\)-cochains} of \(G\) with coefficients in \(M\) is \(C^n(G, M)\). The \emph{\(n\)-coboundary map}
    \begin{align*}
        \dd^n : C^{n-1}(G, M) &\to C^n(G, M)
    \end{align*}
    is dual to \(\dd_n\) in the bar resolution. For \(\phi \in C^{n-1}(G, M)\),
    \begin{align*}
        (\dd^n\phi)(g_1, \dots, g_n) &= g_1\phi(g_2, \dots, g_n) \\
        &- \phi(g_1g_2, \dots, g_n) \\
        &+\phi(g_1, g_2g_3, \dots, g_n) \\
        &+ \dots \\
        &+(-1)^{n-1}\phi(g_1, g_2, \dots, g_{n-1}g_n) \\
        &+(-1)^n\phi(g_1, g_2, \dots, g_{n-1})
    \end{align*}
    The group of \emph{\(n\)-cocycles} is \(Z^n(G, M) = \ker(\dd^{n+1}) \sse C^n(G, M)\), and the \emph{\(n\)-coboundaries} is \(B^n(G, M) = \im(\dd^n) \sse C^n(G, M)\).
\end{definition}
Thus,
\[\HH^n(G, M) = \frac{Z^n(G, M)}{B^n(G, M)}\]

\begin{corollary}
    For any group \(G\), \(\HH^0(G, M) = M^G\) is the invariants.
\end{corollary}

\begin{definition}
    A \emph{derivation of \(G\) with coefficients in \(M\)}  is a function \(\phi : G \to M\), such that
    \[\phi(gh) = g\phi(h) + \phi(g)\]
\end{definition}
Note that \(Z^1(G, M)\) is the set of derivations. Also note that these are also called `crossed homomorphisms'.

An \emph{inner derivation} is \(\phi\) of the form
\[\phi(g) = gm - m\]
for a fixed \(m \in M\).

\begin{corollary}
    \[\HH^1(G, M) = \frac{\{\text{derivations }G \to M\}}{\{\text{inner derivations} G \to M\}}\]
    In particular, if \(M\) is a trivial \(\ZG\)-module, then
    \[\HH^1(G, M) = \Hom_\Group(G, M)\]
\end{corollary}

We will return to considering homology arising from different resolutions.

\begin{definition}
    Let \((A_n, \alpha_n), (B_n, \beta_n)\) be chain complexes of \(\ZG\)-modules, a \emph{chain map} \(f = (f_n)\) are \(\ZG\)-maps \(f_n : A_n \to B_n\), such that
    % https://q.uiver.app/#q=WzAsNCxbMCwwLCJBX24iXSxbMSwwLCJBX3tuLTF9Il0sWzAsMSwiQl9uIl0sWzEsMSwiQl97bi0xfSJdLFswLDEsIlxcYWxwaGFfbiJdLFsyLDMsIlxcYmV0YV9uIiwyXSxbMCwyLCJmX24iLDJdLFsxLDMsImZfe24tMX0iXV0=
\[\begin{tikzcd}
	{A_n} & {A_{n-1}} \\
	{B_n} & {B_{n-1}}
	\arrow["{\alpha_n}", from=1-1, to=1-2]
	\arrow["{\beta_n}"', from=2-1, to=2-2]
	\arrow["{f_n}"', from=1-1, to=2-1]
	\arrow["{f_{n-1}}", from=1-2, to=2-2]
\end{tikzcd}\]
    commutes.
\end{definition}

\begin{lemma}
    Given a chain map \((f_n)\), it induces a map on homology groups
    \[f_* : \HH_n(A_\bullet) \to \HH_n(B_\bullet)\]
\end{lemma}

\lec

\begin{proof}
    Let \(x \in \ker(\alpha_n)\). Define \(f_*([x]) = [f_n(x)]\), where \([\cdot]\) denotes a homology class. Observe that
    \[f_n(x) \in \ker(\beta_n)\]
    since \(\beta_n(f_n(x)) = f_{n-1}(\alpha_n(x)) = 0\). Moreover, if \(x' = x + \alpha_{n-1}(y)\), then
    \[f_n(x') = f_n(x) + f_n(\alpha_{n-1}(y)) = f_n(x) + \beta_n(f_n(y))\]
    and so \([f_n(x')] = [f_n(x)]\). Moreover, this gives a map of abelian groups. 
\end{proof}

\begin{theorem}
    The definition of \(\HH^n(G, M)\) is independent of the choice of projective resolution.
\end{theorem}

\begin{proof}
    Take projective resolutions \((P_n, \dd_n)\) and \((P_n', \dd_n')\) of \(\Z\) by \(\ZG\)-modules. We will produce 
    \begin{itemize}
        \item chain maps \(f_n : P_n \to P_n'\), \(g_n : P_n' \to P_n\),
        \item chain homotopies \(s_n : P_n \to P_{n+1}\), \(s_n' : P_n' \to P_{n+1}'\),
    \end{itemize}
    That is,
    \[\dd s + s \dd = gf - \id \qqtext{and}\dd s' + s'\dd = fg - \id\]
    Using this data, the \(f_n\) define chain maps
    \[f^* : \Hom_G(P_n', M) \to \Hom_G(P_n, M) \qqtext{and} g^* : \Hom_G(P_n', M) \to \Hom_G(P_n, M)\]
    which induce maps between the respective homology groups. Now observe if \(\phi \in \ker(\dd^{n+1}) \in \Hom(P_n, M)\), then
    \begin{align*}
        f^*g^*(\phi)(x) &= \phi(g(f(x))) \\
        &= \phi(x) + \phi(\dd s(x)) + \phi(s \dd (x)) \\
        &= \phi(x) + s^*\dd\phi(x) + \dd s^*\phi(x) \\
        &= \phi(x) + \dd s^*\phi(x)
    \end{align*}
    and so \(f^*g^*(\phi) = \phi + \dd(s^*\phi)\). Hence \(f^*g^*\) induces the identity map on homology. Similarly, \(g^*f^*\) also induce the identity map. This, \(f^*, g^*\) yield isomorphisms on homology. It remains to construct the maps as above.

    Consider the end of the resolutions, so we have
    \(f_{-1} : \Z \to \Z\)
    which is the identity map, and \(f_{-2} : 0 \to 0\) is the zero map. Suppose we have already defined \(f_{n-1}, f_n\). We would like to construct \(f_{n+1}\). Thus, \(f_n \dd : P_{n+1} \to P_n'\), and \(\dd'(f\dd) = f \dd \dd = 0\). Thus, the map \(f\dd\) has image contained in \(\ker(\dd')\). So we can define \(f_{n+1}\) as follows:
    % https://q.uiver.app/#q=WzAsNyxbMSwwLCJQX3tuKzF9Il0sWzIsMCwiUF9uIl0sWzMsMCwiUF97bi0xfSJdLFsyLDEsIlBfbiciXSxbMywxLCJQX3tuLTF9Il0sWzEsMSwiXFxrZXIoXFxkZCcpIl0sWzAsMSwiUF97bisxfSciXSxbMCwxLCJcXGRkIl0sWzEsMiwiXFxkZCJdLFsxLDMsImZfbiJdLFsyLDQsImZfe24tMX0iXSxbMyw0LCJcXGRkJyIsMl0sWzUsMywiIiwyLHsic3R5bGUiOnsidGFpbCI6eyJuYW1lIjoiaG9vayIsInNpZGUiOiJ0b3AifX19XSxbNiw1LCIiLDIseyJzdHlsZSI6eyJoZWFkIjp7Im5hbWUiOiJlcGkifX19XSxbMCw1LCJmX25cXGRkIl0sWzAsNiwiZl97bisxfSIsMix7InN0eWxlIjp7ImJvZHkiOnsibmFtZSI6ImRhc2hlZCJ9fX1dXQ==
    \[\begin{tikzcd}
        & {P_{n+1}} & {P_n} & {P_{n-1}} \\
        {P_{n+1}'} & {\ker(\dd')} & {P_n'} & {P_{n-1}}
        \arrow["\dd", from=1-2, to=1-3]
        \arrow["\dd", from=1-3, to=1-4]
        \arrow["{f_n}", from=1-3, to=2-3]
        \arrow["{f_{n-1}}", from=1-4, to=2-4]
        \arrow["{\dd'}"', from=2-3, to=2-4]
        \arrow[hook, from=2-2, to=2-3]
        \arrow[two heads, from=2-1, to=2-2]
        \arrow["{f_n\dd}", from=1-2, to=2-2]
        \arrow["{f_{n+1}}"', dashed, from=1-2, to=2-1]
    \end{tikzcd}\]
    Here, by exactness the map \(P_{n+1}' \to \ker(\dd')\) is surjective, and so as \(P_{n+1}'\) is projective, such a map exists. We can define \(g_n\) similarly.

    To define \(s_n\), set \(h_n = g_n f_n - \id\). This is a map from \(P_n\) to itself. In particular, \(h_n\) is a chain map, with \(h_{-1} = 0\). Set \(s_{-1} : \Z \to P_0\) to be the zero map, and note \(\dd_0 h_0 = h_{-1}\dd_0 =0\), and so \(\im(h_0) \sse \ker(d_0)\). As before, \(\dd_1 : P_1 \to \ker(\dd_0)\) is surjective, and so we have
    % https://q.uiver.app/#q=WzAsNixbMiwwLCJQXzAiXSxbMiwyLCJQXzAiXSxbMywwLCJcXFoiXSxbMywyLCJcXFoiXSxbMSwyLCJcXGtlcihcXGRkKSJdLFswLDIsIlBfMSJdLFswLDEsImhfMCJdLFsyLDMsIjAiXSxbMCwyXSxbMSwzXSxbNCwxXSxbMCw0LCJoXzAiLDFdLFs1LDQsIiIsMSx7InN0eWxlIjp7ImhlYWQiOnsibmFtZSI6ImVwaSJ9fX1dLFswLDUsInNfMCIsMSx7InN0eWxlIjp7ImJvZHkiOnsibmFtZSI6ImRhc2hlZCJ9fX1dXQ==
\[\begin{tikzcd}
	&& {P_0} & \Z \\
	\\
	{P_1} & {\ker(\dd)} & {P_0} & \Z
	\arrow["{h_0}", from=1-3, to=3-3]
	\arrow["0", from=1-4, to=3-4]
	\arrow[from=1-3, to=1-4]
	\arrow[from=3-3, to=3-4]
	\arrow[from=3-2, to=3-3]
	\arrow["{h_0}"{description}, from=1-3, to=3-2]
	\arrow[two heads, from=3-1, to=3-2]
	\arrow["{s_0}"{description}, dashed, from=1-3, to=3-1]
\end{tikzcd}\]
    Inductively, suppose \(s_{n-1}, s_{n-2}\) are already defined, set \(t_n = h_n - s_{n-1}\dd_n : P_n \to P_n\). We have that 
    \begin{align*}
        \dd_n t_n &= \dd_n h_n - \dd_n s_{n-1}\dd_n \\
        &= h_{n-1}\dd_n - (h_{n-1} - s_{n-2}\dd_{n-1})\dd_n \\
        &= s_{n-2}\dd_{n-1}\dd_n\\
        &= 0
    \end{align*}
    So \(\im(t_n) \sse \ker(\dd_n)\). So we have\footnote{THIS DIAGRAM DOES NOT COMMUTE!}
    % https://q.uiver.app/#q=WzAsNixbMSwwLCJQX24iXSxbMywwLCJQX3tuLTF9Il0sWzIsMiwiUF9uIl0sWzQsMiwiUF97bi0xfSJdLFsxLDIsIlxca2VyKFxcZGRfbikiXSxbMCwyLCJQX3tuKzF9Il0sWzAsMSwiXFxkZF9uIl0sWzIsMywiXFxkZF9uIl0sWzUsNCwiIiwwLHsic3R5bGUiOnsiaGVhZCI6eyJuYW1lIjoiZXBpIn19fV0sWzQsMiwiIiwwLHsic3R5bGUiOnsidGFpbCI6eyJuYW1lIjoiaG9vayIsInNpZGUiOiJ0b3AifX19XSxbMCwyLCJoX24iXSxbMSwzLCJoX3tuLTF9IiwxXSxbMSwyLCJzX3tuLTF9IiwxXSxbMCw0LCJ0X24iLDFdLFswLDUsInNfbiIsMSx7InN0eWxlIjp7ImJvZHkiOnsibmFtZSI6ImRhc2hlZCJ9fX1dXQ==
\[\begin{tikzcd}
	& {P_n} && {P_{n-1}} \\
	\\
	{P_{n+1}} & {\ker(\dd_n)} & {P_n} && {P_{n-1}}
	\arrow["{\dd_n}", from=1-2, to=1-4]
	\arrow["{\dd_n}", from=3-3, to=3-5]
	\arrow[two heads, from=3-1, to=3-2]
	\arrow[hook, from=3-2, to=3-3]
	\arrow["{h_n}", from=1-2, to=3-3]
	\arrow["{h_{n-1}}"{description}, from=1-4, to=3-5]
	\arrow["{s_{n-1}}"{description}, from=1-4, to=3-3]
	\arrow["{t_n}"{description}, from=1-2, to=3-2]
	\arrow["{s_n}"{description}, dashed, from=1-2, to=3-1]
\end{tikzcd}\]
    Define \(s_n'\) similarly and we are done.
\end{proof}

\begin{remark}
    For any (left) \(\ZG\)-module \(N\), we can take a projective (resp. free) resolution of \(N\) by \(\ZG\)-modules. Repeating everything we have done, and applying \(\Hom_G(\cdot, M)\) gives homology groups, which are called
    \[\Ext^n_{\ZG}(N, M)\]
    and so
    \[\HH^n(G, M) = \Ext^n_\ZG(\Z, M)\]
    As above, the \(\Ext\)-groups are independent of resolution.
\end{remark}

\section{Low degree cohomology, group extensions}

\subsection{First cohomology}

Recall that \(\HH^0(G, M) = M^G\) is the group of invariants under \(G\), and that a derivation (or a \(1\)-cocycle) is a map \(\phi : G \to M\), such that
\[\phi(g_1g_2) = g_1\phi(g_2) + \phi(g_1)\]
We will see two interpretations of (inner) derivations. Recall that an inner derivation is one of the form \(\phi(g) = gm - m\), for some fixed \(m \in M\).

Let \(M\) be a \(\ZG\)-module, and consider possible \(\ZG\)-actions on the abelian group \(M \oplus \Z\) of the form
\[g(m, n) = (gm + n\phi(g), n)\]
In this case,
\begin{align*}
    g_1g_2(m, n) = g_1(g_2m + n\phi(g_2), n) = (g_1g_2m + ng_1\phi(g_2) + n\phi(g_1), n)
\end{align*}
On the other hand,
\[(g_1g_2)(m, n) = (g_1g_2 m + n\phi(g_1g_2), n)\]
These are the same exactly when \(\phi\) is a derivation. In particular, if \(M\) is a free \(\Z\)-module of finite rank, then we get a map \(G \to \GL(M)\), by can write this as
\[g \mapsto \begin{pmatrix}
    \theta_1(g) & \dd g \\
    0 & 1
\end{pmatrix}\]
where \(\theta_1(g)\) is the action of \(g\) on \(M\). In particular, this is a homomorphism of groups if and only if \(\phi\) is a derivation. Moreover, \(\phi\) is an inner derivation if and only if \((-m, 1)\) generates a \(\ZG\)-submodule which is a trivial module inside \(M \oplus \Z\).

\lec

For the second interpretation:

\begin{definition}
    [semidirect product]
    Let \(G\) be a group, \(M\) be a \(\ZG\)-module. We can construct their \emph{semidirect product} \(M \rtimes G\) as a group \(M \times G\) with operation
    \[(m_1, g_1) * (m_2, g_2) = (m_1 + g_1m_2, g_1g_2)\]
\end{definition}
Here, we have an isomorphism 
\[M \cong \{(m, 1) \mid m \in M\}\]
This is an abelian normal subgroup. Similarly, we have an isomorphism
\[G \cong \{(0, g) \mid g \in G\}\]
and conjugation by \(\{(0, g)\}\) corresponds to the \(G\)-action on \(M\).

Moreover,
\[\frac{M \rtimes G}{\{(m, 1) \mid m \in M\}} \cong G\]
Note that there is a group homomorphism \(s : G \to M \rtimes G\), \(s(g) = (0, g)\) is such that
\[G \to M \rtimes G \to G\]
is the identity. \(s\) is called a \emph{splitting}. Now if we have another splitting \(s_1 : G \to M \rtimes G\), so that the composition \(G \to M \rtimes G \to G\) is the identity. Define \(\psi_{s_1} : G \to M\), so that
\[s_1(g) = (\psi_{s_1}(g), g)\]
Then \(\psi_{s_1} \in Z^1(G, M)\), and given two splittings \(s_1, s_2\), \(\psi_{s_1} - \psi_{s_2} \in B^1(G, M)\) if and only if there exists an \(m\) such that
\[(m, 1)s_1(g)(m_1)^{-1} = s_2(g)\]
Conversely, given \(\phi \in Z^1(G, M)\), there exists a splitting \(s_1\) such that \(\phi = \psi_{s_1}\).

\begin{theorem}
    \(\HH^1(G, M)\) correspond to the \(M\)-conjugacy classes of splittings.
\end{theorem}

\begin{proof}
    Examples Sheet 1.
\end{proof}

\subsection{Second cohomology}

\begin{definition}
    [extension] An \emph{extension} \(E\) of \(G\) by \(M\) for a group \(G\), and a \(\ZG\)-module \(M\), is an exact sequence
    % https://q.uiver.app/#q=WzAsNSxbMCwwLCIwIl0sWzEsMCwiTSJdLFsyLDAsIkUiXSxbMywwLCJHIl0sWzQsMCwiMSJdLFswLDFdLFsxLDJdLFsyLDNdLFszLDRdXQ==
\[\begin{tikzcd}
	0 & M & E & G & 1
	\arrow[from=1-1, to=1-2]
	\arrow[from=1-2, to=1-3]
	\arrow[from=1-3, to=1-4]
	\arrow[from=1-4, to=1-5]
\end{tikzcd}\]
    where the maps are group homomorphisms. That is, \(M\) embeds in \(E\), so that the image \(M\) is an abelian normal subgroup. \(E\) acts on \(M\) by conjugation, and as \(M\) is abelian, we have an induced action on \(E/M \cong G\). This agrees with the given \(G\)-action on \(M\).
\end{definition}

\begin{example}
    \(E = M \rtimes G\) is an extension of \(G\) by \(M\). This is called the \emph{split extension}.
\end{example}

\begin{definition}
    [equivalent]
    Two extensions are \emph{equivalent} if we have a commuting diagram of group homomorphisms
    % https://q.uiver.app/#q=WzAsNixbMCwxLCIwIl0sWzEsMSwiTSJdLFsyLDAsIkUiXSxbMywxLCJHIl0sWzQsMSwiMSJdLFsyLDIsIkUnIl0sWzAsMV0sWzEsMl0sWzIsM10sWzMsNF0sWzEsNV0sWzUsM10sWzIsNSwiIiwxLHsiY29sb3VyIjpbMCw2MCw2MF19XV0=
    \[\begin{tikzcd}
        && E \\
        0 & M && G & 1 \\
        && {E'}
        \arrow[from=2-1, to=2-2]
        \arrow[from=2-2, to=1-3]
        \arrow[from=1-3, to=2-4]
        \arrow[from=2-4, to=2-5]
        \arrow[from=2-2, to=3-3]
        \arrow[from=3-3, to=2-4]
        \arrow[color={rgb,255:red,214;green,92;blue,92}, from=1-3, to=3-3]
    \end{tikzcd}\] 
\end{definition}

\begin{remark}
    The vertical map is an isomorphism of groups. However, the converse is false. On examples sheet 1, we have \(E, E'\) isomorphic groups, but inequivalent extensions.
\end{remark}

\begin{definition}
    [central extension] A \emph{central} extension is one where the \(\ZG\)-module is a trivial module.
\end{definition}

\begin{proposition}
    Let \(E\) be an extension of \(G\) by \(M\). If there exists a splitting \(s_1 : G \to E\), then \(E\) is equivalent to \(M \rtimes G\).
\end{proposition}

\begin{proof}
    Exercise.
\end{proof}

\begin{theorem}
    Let \(G\) be a group, and \(M\) be a \(\ZG\)-module. Then there exists a bijection between \(\HH^2(G, M)\) and the equivalence classes of extensions of \(G\) by \(M\).
\end{theorem}

Given an extenstion
    % https://q.uiver.app/#q=WzAsNSxbMCwwLCIwIl0sWzEsMCwiTSJdLFsyLDAsIkUiXSxbMywwLCJHIl0sWzQsMCwiMSJdLFswLDFdLFsxLDJdLFsyLDNdLFszLDRdXQ==
    \[\begin{tikzcd}
        0 & M & E & G & 1
        \arrow[from=1-1, to=1-2]
        \arrow[from=1-2, to=1-3]
        \arrow[from=1-3, to=1-4]
        \arrow[from=1-4, to=1-5]
    \end{tikzcd}\]
a \emph{section} is a function \(s : G \to E\), such that the composition \(\pi \circ s = \id\), where \(\pi\) is the map \(E \to G\). Note that \(s\) need not be a group homomorphism. Suppose \(s(1) = 1\). Let
\[\phi(g_1,g_2) = s(g_1)s(g_2)s(g_1g_2)^{-1}\]
Then
\[\pi(\phi(g_1g_2)) = 1\]
and so \(\phi(g_1, g_2) \in \ker(\pi) = M\). That is, we have a map \(\phi : G^2 \to M\), which is a \(2\)-cochain. In fact, it is a \(2\)-cocycle.
\begin{align*}
    s(g_1)s(g_2)s(g_2) &= \phi(g_1, g_2)s(g_1g_2)s(g_3) \\
    &= \phi(g_1,g_2)\phi(g_1g_2,g_3)s(g_1g_2g_3) \tag{\(*\)}
\end{align*}
Similarly,
\begin{align*}
    s(g_1)s(g_2)s(g_3) &= s(g_1)\phi(g_2, g_2)s(g_2g_3) \\
    &= s(g_1)\phi(g_2g_3)s(g_1)^{-1}s(g_1)s(g_2g_3) \\
    &= s(g_1)\phi(g_2g_2)s(g_1)^{-1}\phi(g_1, g_2g_3)s(g_1g_2g_3) \tag{\(**\)}
\end{align*}
Equating \((*)\) and \((**)\), cancelling the \(s(g_1g_2g_3)\), and changing into additive notation, we get
\[\phi(g_1, g_2) + \phi(g_1g_2, g_3) = g_1\phi(g_2, g_3) + \phi(g_1, g_2g_3)\]
Hence
\[\dd\phi(g_1, g_2, g_3) = 0\]
and so \(\phi \in Z^2(G, M)\). Note \(\phi\) is \emph{normalised}. That is, \(\phi(1, g) = \phi(g, 1) = 0\). So what we have shown is that an extension of \(G\) by \(M\), along with a choice of section \(s\), yields a normalised 2-cocycle \(\phi \in Z^2(G, M)\).

Now take another choice of section \(s'\), with \(s'(1) = 1\). Then the corresponding normalised cocycles \(\phi\) and \(\phi'\) differ by a coboundary, and so we have defined a map
\[\{\text{extensions}\} \to \HH^2(G, M)\]

\lec

To see this, note that \(\pi(s(g)s'(g)^{-1}) = 1\), and \(\psi(g) = s(g)s'(g)^{-1} \in \ker(\pi) = M\). Thus, we have a map \(\psi : G \to M\). Then
\begin{align*}
    s'(g_1)s'(g_2) &= \psi(g_1)s(g_1)\psi(g_2)s(g_2) \\
    &= \psi(g_1)s(g_1)\psi(g_2)s(g_1)^{-1}s(g_1)s(g_2) \\
    &= \psi(g_1)s(g_1)\psi(g_2)s(g_1)^{-1}\phi(g_1, g_2)s(g_1g_2) \\
    &= \psi(g_1)s(g_1)\psi(g_2)s(g_2)^{-1}\phi(g_1, g_2)\psi(g_1g_2)^{-1}s'(g_1g_2)
\end{align*}
Switching to additive notation:
\begin{align*}
    \phi'(g_1, g_2) &= \psi(g_1) + g_1\psi(g_2) + \phi(g_1, g_2) - \psi(g_1g_2) \\
    &= \phi(g_1, g_2) + \dd\psi(g_1, g_2)
\end{align*}
Thus, \(\phi\) and \(\phi'\) differ by a coboundary. For the rest of the proof, we need:
\begin{enumerate}[(a)]
    \item to show that equivalent extensions give the same cohomology class,
    \item construct an inverse map from cohomology classes to extensions,
    \item show that the two maps are inverse to each other.
\end{enumerate}
We will show (b), and leave (a) and (c). For this, we need
\begin{lemma}
    Let \(\phi \in Z^2(G, M)\). Then there exists a cochain \(\psi \in C^1(G, M)\), such that \(\phi + \dd\psi\) is a normalised cocycle. Hence every cohomology class can be reresented by a normalised cocycle.
\end{lemma}
\begin{proof}
    Let \(\psi(g) = -\phi(1, g)\). Then
    \[(\phi + \dd\psi)(1, g) = \phi(1, g) - \left(\phi(1, g) - \phi(1, g) + \phi(1, 1)\right) = \phi(1, g) - \phi(1, 1)\]
    and
    \[(\phi + \dd\psi)(g, 1) = \phi(g, 1) - g\phi(1, 1)\]
    But we know that \(\dd\phi(1, 1, g) = \dd\phi(g, 1, 1) = 0\) as \(\phi\) is a cocycle. A computation shows that both of the above are zero.
\end{proof}

Now take a normalised cocycle \(\phi\) representating our cohomology class. We construct an extension as follows.
% https://q.uiver.app/#q=WzAsNSxbMCwwLCIwIl0sWzEsMCwiTSJdLFsyLDAsIkVfXFxwaGkiXSxbMywwLCJHIl0sWzQsMCwiMSJdLFswLDFdLFsxLDJdLFsyLDNdLFszLDRdXQ==
\[\begin{tikzcd}
	0 & M & {E_\phi} & G & 1
	\arrow[from=1-1, to=1-2]
	\arrow[from=1-2, to=1-3]
	\arrow[from=1-3, to=1-4]
	\arrow[from=1-4, to=1-5]
\end{tikzcd}\]
by considering \(E = M \times G\), with the product structure
\[(m_1, g_1)(m_2, g_2) = (m_1 + g_1m_2 + \phi(g_1, g_2), g_1g_2)\]
For this to be a group operation, we use that \(\phi\) is normalised. In this case, we do have an extension, where the map \(\pi\) is projection onto the second factor. Notice if we have another normalised \(2\)-cocycle \(\phi'\) representing our cohomology class, then \(\phi - \phi' = \dd\psi\) is a coboundary, and we can define a map
\begin{align*}
    E_\phi &\to E_{\phi'} \\
    (m, g) &\mapsto (m + \psi(g), g)
\end{align*}
This gives us an equivalence of extensions.

\begin{example}
    [central extensions of \(\Z^2\) by \(\Z\)] Note we already know of (at least) two such. We could have
    % https://q.uiver.app/#q=WzAsNSxbMCwwLCIwIl0sWzEsMCwiXFxaIl0sWzIsMCwiXFxaIFxcb3BsdXMgXFxaXjIiXSxbMywwLCJcXFpeMiJdLFs0LDAsIjAiXSxbMCwxXSxbMSwyXSxbMiwzXSxbMyw0XV0=
    \[\begin{tikzcd}
        0 & \Z & {\Z \oplus \Z^2} & {\Z^2} & 0
        \arrow[from=1-1, to=1-2]
        \arrow[from=1-2, to=1-3]
        \arrow[from=1-3, to=1-4]
        \arrow[from=1-4, to=1-5]
    \end{tikzcd}\]
    and we also have the Heisenberg group
    \[H = \left\{\begin{pmatrix}
        1 & r & m \\
        & 1 & s \\
        & & 1
    \end{pmatrix}\ \bigg\vert\ r, s, m \in \Z\right\}\]
    and we have an extension given by
    \begin{align*}
        m &\mapsto \begin{pmatrix}
            1 & 0 & m \\
            & 1 & 0 \\
            & & 1
        \end{pmatrix} \\
        \begin{pmatrix}
            1 & r & 0 \\
            & 1 & 0 \\
            & & 1
        \end{pmatrix} &\mapsto (r, s)
    \end{align*}
    Now if we write things multiplicatively, \(T \cong \Z^2\) is generated by \(a, b\). We have a free resolution
    % https://q.uiver.app/#q=WzAsNixbMCwwLCIwIl0sWzEsMCwiXFxaIFQiXSxbMiwwLCIoXFxaIFQpXjIiXSxbMywwLCJcXFogVCJdLFs0LDAsIlxcWiJdLFs1LDAsIjAiXSxbMCwxXSxbMSwyLCJcXGJldGEiXSxbMiwzLCJcXGFscGhhIl0sWzMsNCwiXFxlcHNpbG9uIl0sWzQsNV1d
\[\begin{tikzcd}
	0 & {\Z T} & {(\Z T)^2} & {\Z T} & \Z & 0
	\arrow[from=1-1, to=1-2]
	\arrow["\beta", from=1-2, to=1-3]
	\arrow["\alpha", from=1-3, to=1-4]
	\arrow["\epsilon", from=1-4, to=1-5]
	\arrow[from=1-5, to=1-6]
\end{tikzcd}\]
    where
    \begin{align*}
        \beta(z) &= (z(1 - b), z(a - 1)) \\
        \alpha(x, y) &= x(a - 1) + y(b - 1)
    \end{align*}
    and \(\epsilon\) is the augmentation map. Now applying \(\Hom_T(\cdot, \Z)\), we get the chain complex
    % https://q.uiver.app/#q=WzAsNCxbMCwwLCJcXEhvbV9UKFxcWiBULCBcXFopIl0sWzEsMCwiXFxIb21fVCgoXFxaIFQpXjIsIFxcWikiXSxbMiwwLCJcXEhvbV9UKFxcWiBULCBcXFopIl0sWzMsMCwiMCJdLFswLDEsIlxcYWxwaGFeKiJdLFsxLDIsIlxcYmV0YV4qIl0sWzIsM11d
    \[\begin{tikzcd}
        {\Hom_T(\Z T, \Z)} & {\Hom_T((\Z T)^2, \Z)} & {\Hom_T(\Z T, \Z)} & 0
        \arrow["{\alpha^*}", from=1-1, to=1-2]
        \arrow["{\beta^*}", from=1-2, to=1-3]
        \arrow[from=1-3, to=1-4]
    \end{tikzcd}\]
    In fact, \(\alpha^*, \beta^*\) are the zero maps, and so \(\HH^2(T, \Z) = \Hom_T(\Z T, \Z) \cong \Z\), with generator corresponding to the augmentation map \(\epsilon : \Z T \to \Z\).

    To see that \(\alpha^*, \beta^*\) are zero, take a \(\Z T\)-map \(f : (\Z T)^2 \to \Z\) and \(z \in \Z T\), then
    \begin{align*}
        (\beta^*f)(z) &= f(\beta(z)) \\
        &= f(z(1 - b), z( a - 1)) \\
        &= f(z - bz, 0) + f(0, za - z) \\
        &= (1 - b)f(z, 0) + (a - 1)f(0, z) = 0
    \end{align*}
    The proof that \(\alpha^*\) is zero is similar.

    Next, we will interpret \(\HH^2(T, \Z)\) in terms of cocycles from the bar resolution. We will construct a chain map
    % https://q.uiver.app/#q=WzAsMTAsWzAsMCwiXFxaIFRcXHtUXnsoMil9XFx9Il0sWzEsMCwiXFxaIFRcXHtUXnsoMSl9XFx9Il0sWzIsMCwiXFxaIFRcXHtUXnsoMCl9XFx9Il0sWzMsMCwiXFxaIl0sWzQsMCwiMCJdLFswLDEsIlxcWiBUIl0sWzEsMSwiKFxcWiBUKV4yIl0sWzIsMSwiXFxaIFQiXSxbMywxLCJcXFoiXSxbNCwxLCIwIl0sWzAsMV0sWzEsMl0sWzIsMywiXFxlcHNpbG9uIl0sWzMsNF0sWzUsNiwiXFxiZXRhIiwyXSxbNiw3LCJcXGFscGhhIiwyXSxbMiw3LCJcXGlkIiwyXSxbNyw4XSxbMyw4LCI9IiwyXSxbOCw5XSxbMSw2LCJmXzEiLDJdLFswLDUsImZfMiIsMl1d
\[\begin{tikzcd}
	{\Z T\{T^{(2)}\}} & {\Z T\{T^{(1)}\}} & {\Z T\{T^{(0)}\}} & \Z & 0 \\
	{\Z T} & {(\Z T)^2} & {\Z T} & \Z & 0
	\arrow[from=1-1, to=1-2]
	\arrow[from=1-2, to=1-3]
	\arrow["\epsilon", from=1-3, to=1-4]
	\arrow[from=1-4, to=1-5]
	\arrow["\beta"', from=2-1, to=2-2]
	\arrow["\alpha"', from=2-2, to=2-3]
	\arrow["\id"', from=1-3, to=2-3]
	\arrow[from=2-3, to=2-4]
	\arrow["{=}"', from=1-4, to=2-4]
	\arrow[from=2-4, to=2-5]
	\arrow["{f_1}"', from=1-2, to=2-2]
	\arrow["{f_2}"', from=1-1, to=2-1]
\end{tikzcd}\]
    Note that for degrees \(-1\) and \(0\), we can just take the identity map. To construct \(f_1 : \Z T\{T^{(1)}\} \to (\Z T)^2\), we need that \(\alpha f_1 = \dd\). So we just need to define the image of the symbols \([a^rb^s]\), for \(r, s \in \Z\). Say
    \[f_1([a^ry^s]) = (x_{r,s}, y_{r, s})\]
    and we need
    \[\alpha(x_{r, s}, y_{r, s}) = \dd[a^rb^s] = a^rb^s - 1 = (a^r - 1)b^s + (b^s - 1)\]
    Define
    \begin{align*}
        S(a, r) = \begin{cases}
            1 + a + \cdots + a^{r - 1} & r > 0 \\
            -a^{-1} - \cdots - a^r & r \le 0 \\
        \end{cases}
    \end{align*}
    We note that \(S(a, r)(a - 1) = a^r - 1\) in both cases. Then
    \[\alpha(S(a, r)b^s, S(b, s)) = S(a, r)b^s(a - 1) + S(b, s)(b - 1) = \dd[a^rb^s]\]
    as required. So we can define
    \[f_1[a^rb^s] = (S(a, r)b^s, S(b, s))\]
\end{example}

\printindex

\end{document}